
\pdfminorversion=4 % for acroread
\ifdefined\ishandout
  \documentclass[aspectratio=169,t,handout,xcolor={usenames,dvipsnames}]{beamer}
\else
  \documentclass[aspectratio=169,t,xcolor={usenames,dvipsnames}]{beamer}
\fi

\usepackage{../beamerstyle}
\usepackage{dsfont}
\usepackage{bm}
\usepackage[english]{babel}
\usepackage[utf8]{inputenc}
\usepackage{graphicx}
\usepackage{algorithm}
\usepackage[ruled,vlined,algo2e,linesnumbered]{algorithm2e}
%\usepackage[boxed,vlined]{algorithm2e}
\usepackage{hyperref}
\usepackage{booktabs}
\usepackage{mathtools}

\usepackage{amsmath,amssymb}
\usepackage{listings}
\lstset{frame=lines,framesep=3pt,numbers=left,numberblanklines=false,basicstyle=\ttfamily\small}

\usepackage{subfig}
\usepackage{multicol}
%\usepackage{appendixnumberbeamer}
%
\usepackage{tcolorbox}

\usepackage{pgfplots}
\usepackage{tikz}
\usetikzlibrary{trees} 
\usetikzlibrary{shapes.geometric}
\usetikzlibrary{positioning,shapes,shadows,arrows,calc,mindmap}
\usetikzlibrary{positioning,fadings,through}
\usetikzlibrary{decorations.pathreplacing}
\usetikzlibrary{intersections}
\usetikzlibrary{positioning,fit,calc,shadows,backgrounds}
\pgfdeclarelayer{background}
\pgfdeclarelayer{foreground}
\pgfsetlayers{background,main,foreground}
\tikzstyle{activity}=[rectangle, draw=black, rounded corners, text centered, text width=8em]
\tikzstyle{data}=[rectangle, draw=black, text centered, text width=8em]
\tikzstyle{myarrow}=[->, thick, draw=black]

% Define the layers to draw the diagram
\pgfdeclarelayer{background}
\pgfdeclarelayer{foreground}
\pgfsetlayers{background,main,foreground}

%\usepackage{listings}
%\lstset{numbers=left,
%  showstringspaces=false,
%  frame={tb},
%  captionpos=b,
%  lineskip=0pt,
%  basicstyle=\ttfamily,
%%  extendedchars=true,
%  stepnumber=1,
%  numberstyle=\small,
%  xleftmargin=1em,
%  breaklines
%}

 
\definecolor{blue}{RGB}{0, 74, 153}

\usetheme{Boadilla}
%\useinnertheme{rectangles}
\usecolortheme{whale}
\setbeamercolor{alerted text}{fg=blue}
\useoutertheme{infolines}
\setbeamertemplate{navigation symbols}{\vspace{-5pt}} % to lower the logo
\setbeamercolor{date in head/foot}{bg=white} % blue
\setbeamercolor{date in head/foot}{fg=blue}
\setbeamercolor{author  in head/foot}{bg=white} %blue
\setbeamercolor{title in head/foot}{bg=white} % blue
\setbeamercolor{title}{fg=white, bg=blue}
\setbeamercolor{block title}{fg=white,bg=blue}
\setbeamercolor{block body}{bg=blue!10}
\setbeamercolor{frametitle}{fg=white, bg=blue}
\setbeamercovered{invisible}

% \makeatletter
% \setbeamertemplate{footline}
% {
%   \leavevmode%
%   \hbox{%
%   \begin{beamercolorbox}[wd=.333333\paperwidth,ht=2.25ex,dp=1ex,center]{author in head/foot}%
% %    \usebeamerfont{author in head/foot}\insertshortauthor
%   \end{beamercolorbox}%
%   \begin{beamercolorbox}[wd=.333333\paperwidth,ht=2.25ex,dp=1ex,center]{title in head/foot}%
%     \usebeamerfont{title in head/foot}\insertshorttitle
%   \end{beamercolorbox}%
%   \begin{beamercolorbox}[wd=.333333\paperwidth,ht=2.25ex,dp=1ex,right]{date in head/foot}%
%     \usebeamerfont{date in head/foot}\insertshortdate{}\hspace*{2em}
% %    \insertframenumber\hspace*{2ex} 
%   \end{beamercolorbox}}%
%   \vskip0pt%
% }
% \makeatother

%\pgfdeclareimage[height=1.2cm]{automl}{images/logos/automl.png}
%\pgfdeclareimage[height=1.2cm]{freiburg}{images/logos/freiburg}

%\logo{\pgfuseimage{freiburg}}

\renewcommand{\comment}[1]{
	\noindent
	%\vspace{0.25cm}
	{\color{red}{\textbf{TODO:} #1}}
	%\vspace{0.25cm}
}
\newcommand{\notefh}[1]{\textcolor{red}{\textbf{FH:} #1}}
\renewcommand{\comment}[1]{}
\newcommand{\hide}[1]{}
\newcommand{\cemph}[2]{\emph{\textcolor{#1}{#2}}}

\newcommand{\lit}[1]{{\footnotesize\color{black!60}[#1]}}

\newcommand{\litw}[1]{{\footnotesize\color{blue!20}[#1]}}


\newcommand{\myframe}[2]{\begin{frame}[c]{#1}#2\end{frame}}
\newcommand{\myframetop}[2]{\begin{frame}{#1}#2\end{frame}}
\newcommand{\myit}[1]{\begin{itemize}#1\end{itemize}}
\newcommand{\myblock}[2]{\begin{block}{#1}#2\end{block}}


\newcommand{\votepurple}[1]{\textcolor{Purple}{$\bigstar$}}
\newcommand{\voteyellow}[1]{\textcolor{Goldenrod}{$\bigstar$}}
\newcommand{\voteblue}[1]{\textcolor{RoyalBlue}{$\bigstar$}}
\newcommand{\votepink}[1]{\textcolor{Pink}{$\bigstar$}}

\newcommand{\diff}{\mathop{}\!\mathrm{d}}
\newcommand{\refstyle}[1]{{\small{\textcolor{gray}{#1}}}}
\newcommand{\hands}[0]{\includegraphics[height=1.5em]{images/hands}}
\newcommand{\transpose}[0]{{\textrm{\tiny{\sf{T}}}}}
\newcommand{\norm}{{\mathcal{N}}}
\newcommand{\cutoff}[0]{\kappa}
\newcommand{\instD}[0]{\dataset}
\newcommand{\insts}[0]{\mathcal{I}}
\newcommand{\inst}[0]{i}
\newcommand{\instI}[1]{i^{(#1)}}

% Iteration specific instance of variable/function/anything
% Introduced in the BO section, but moved up here to make it available within other macros
\newcommand{\iter}[2][\bocount]{{#2}^{(#1)}}

%--------HPO parameter macros-----------

% Parameter Configuration Space
\newcommand{\pcs}[0]{\pmb{\Lambda}}

% ???
\newcommand{\bx}[0]{\conf}

% Parameter Configuration
\newcommand{\conf}[0]{\pmb{\lambda}}

% Final Configuration
\newcommand{\finconf}[0]{\pmb{\hat{\lambda}}}

% Configuration corresponding to a given iteration -- better use \iter!
\newcommand{\confI}[1]{{\conf}^{(#1)}}

% Default Configuration
\newcommand{\defconf}[0]{{\conf}_{\text{def}}}

% Incumbent Configuration
\newcommand{\incumbent}[1][\bocount]{\iter[#1]{\finconf}}

% Optimal Configuration
\newcommand{\optconf}[0]{{\conf}^*}

% Configuration Space
\newcommand{\confs}[0]{\pcs}

%----------------------------------------

%\newcommand{\vlambda}[0]{\bm{\lambda}}
%\newcommand{\vLambda}[0]{\bm{\Lambda}}
\newcommand{\dataset}[0]{\mathcal{D}}
\newcommand{\datasets}[0]{\mathbf{D}}
\newcommand{\loss}[0]{L}
\newcommand{\risk}{\mathcal{R}}
\newcommand{\riske}{\mathcal{R}_{\text{emp}}}
\newcommand{\cost}[0]{c}
\newcommand{\costI}[1]{c^{(#1)}}

% Gaussian Process
\newcommand{\gp}{\mathcal{G}}
% Family of Objective Functions
\newcommand{\objF}{F}

%---------------BO Macros------------------

% BO loop counter
\newcommand{\bocount}{t}
% BO loop counter max, the counter runs from 1 to this value
\newcommand{\bobudget}{T}
% BO loop observation
\newcommand{\obs}[1][\conf]{\cost({#1})}
% BO loop observation space
\newcommand{\obsspace}{\mathcal{Y}}
% BO loop next observation
\newcommand{\bonextobs}{\obs[\iter{\conf}]}
% Acquisition Function, no args
\newcommand{\acq}{u}
% Standard Normal PDF
\newcommand{\pdf}{\phi}
% Standard Normal CDF
\newcommand{\cdf}{\Phi}
% Mean
\newcommand{\mean}{\mu}
% Standard Deviation
\newcommand{\stddev}{\sigma}
% Variance
\newcommand{\variance}{\sigma^2}
% Noise
\newcommand{\noise}{\nu}
% BO loop next selected sample
\newcommand{\bonextsample}{\confI{\bocount}}

% Single hyperparameter
\newcommand{\hyperparam}{\lambda}

% Single hyperparameter within a hyperparameter configuration
\newcommand{\hyperparami}[1][i]{{\hyperparam}_#1}

% Full definition of final configuration
\newcommand{\finconffull}{\incumbent[\bobudget]}

% Dataset
\newcommand{\datasetHPO}{{\dataset}_{HPO}}

% Dataset definition
\newcommand{\datasetHPOdef}{{\langle \bonextsample,\,\bonextobs \rangle}_{\bocount=1}^{\bobudget}}

% Double Display Fraction, forces large displays for everything in numerator and denominator
\newcommand\ddfrac[2]{\frac{\displaystyle #1}{\displaystyle #2}}

% Conditional Probability "Given That" Relation, source:https://tex.stackexchange.com/a/141685/205886
\newcommand\given[1][]{\:#1\vert\:}

% Expectation as a math operator
\DeclareMathOperator*{\E}{\mathbb{E}}

% Citation 
\newcommand{\source}[1]{
    \begin{flushright}
    	Source: \lit{#1}
    \end{flushright}
}
%-------------------------------------------

%Real numbers set
\newcommand{\realnum}{\mathbb{R}}
%Configuration space - do not use
%\newcommand{\configspace}{\Theta}
%Instances - do not use
%\newcommand{\instances}{\mathcal{I}}
%Expected value
\newcommand{\expectation}{\mathbb{E}}
%Kernel
\newcommand{\kernel}{\kappa}
%Constraint function
\newcommand{\constraintf}{c}
%Normal distribution
\newcommand{\normaldist}{\mathcal{N}}

% \renewcommand{\vec}[1]{\mathbf{#1}}
\newcommand{\hist}[0]{\dataset_{\text{Hist}}}
\newcommand{\param}[0]{p}
\newcommand{\algo}[0]{\mathcal{A}}
\newcommand{\algos}[0]{\mathbf{A}}
%\newcommand{\nn}[0]{N}
\newcommand{\feats}[0]{\mathcal{X}_{\text{meta}}}
\newcommand{\feat}[0]{\x_{\text{meta}}}
%\newcommand{\cluster}[0]{\vec{h}}
%\newcommand{\clusters}[0]{\vec{H}}
\newcommand{\perf}[0]{\mathbb{R}}
%\newcommand{\surro}[0]{\mathcal{S}}
\newcommand{\surro}[0]{\hat{\cost}}
\newcommand{\func}[0]{f}
\newcommand{\epm}[0]{\surro}
\newcommand{\portfolio}[0]{\mathbf{P}}
\newcommand{\schedule}[0]{\mathcal{S}}

% Machine Learning
\newcommand{\mdata}[0]{\dataset_{\text{meta}}}
\newcommand{\datasettrain}[0]{\dataset_{\text{train}}}
\newcommand{\datasetval}[0]{\dataset_{\text{val}}}
\newcommand{\datasettest}[0]{\dataset_{\text{test}}}
\newcommand{\x}[0]{\mathbf{x}}
\newcommand{\y}[0]{y}
\newcommand{\xI}[1]{\mathbf{x}^{(#1)}}
\newcommand{\yI}[1]{y^{(#1)}}
\newcommand{\fx}{f(\mathbf{x})}  % f(x), continuous prediction function
\newcommand{\Hspace}{\mathcal{H}} % hypothesis space where f is from
\newcommand{\fh}{\hat{f}}       % f hat, estimated prediction function

% Deep Learning
\newcommand{\weights}[0]{\theta}
\newcommand{\metaweights}[0]{\phi}


% reinforcement learning
\newcommand{\policies}[0]{\mathbf{\Pi}}
\newcommand{\policy}[0]{\pi}
\newcommand{\actionRL}[0]{a}
\newcommand{\stateRL}[0]{s}
\newcommand{\statesRL}[0]{\mathcal{S}}
\newcommand{\rewardRL}[0]{r}
\newcommand{\rewardfuncRL}[0]{\mathcal{R}}

\RestyleAlgo{algoruled}
\DontPrintSemicolon
\LinesNumbered
\SetAlgoVlined
\SetFuncSty{textsc}

\SetKwInOut{Input}{Input}
\SetKwInOut{Output}{Output}
\SetKw{Return}{return}

%\newcommand{\changed}[1]{{\color{red}#1}}

%\newcommand{\citeN}[1]{\citeauthor{#1}~(\citeyear{#1})}

\renewcommand{\vec}[1]{\mathbf{#1}}
\DeclareMathOperator*{\argmin}{arg\,min}
\DeclareMathOperator*{\argmax}{arg\,max}

%\newcommand{\aqme}{\textit{AQME}}
%\newcommand{\aslib}{\textit{ASlib}}
%\newcommand{\llama}{\textit{LLAMA}}
%\newcommand{\satzilla}{\textit{SATzilla}}
%\newcommand{\satzillaY}[1]{\textit{SATzilla'{#1}}}
%\newcommand{\snnap}{\textit{SNNAP}}
%\newcommand{\claspfolioTwo}{\textit{claspfolio~2}}
%\newcommand{\flexfolio}{\textit{FlexFolio}}
%\newcommand{\claspfolioOne}{\textit{claspfolio~1}}
%\newcommand{\isac}{\textit{ISAC}}
%\newcommand{\eisac}{\textit{EISAC}}
%\newcommand{\sss}{\textit{3S}}
%\newcommand{\sunny}{\textit{Sunny}}
%\newcommand{\ssspar}{\textit{3Spar}}
%\newcommand{\cshc}{\textit{CSHC}}
%\newcommand{\cshcpar}{\textit{CSHCpar}}
%\newcommand{\measp}{\textit{ME-ASP}}
%\newcommand{\aspeed}{\textit{aspeed}}
%\newcommand{\autofolio}{\textit{AutoFolio}}
%\newcommand{\cedalion}{\textit{Cedalion}}
\newcommand{\fanova}{\textit{fANOVA}}
\newcommand{\sbs}{\textit{SB}}
\newcommand{\oracle}{\textit{VBS}}

% like approaches
\newcommand{\claspfoliolike}[1]{\texttt{claspfolio-#1-like}}
\newcommand{\satzillalike}[1]{\texttt{SATzilla'#1-like}}
\newcommand{\isaclike}{\texttt{ISAC-like}}
\newcommand{\ssslike}{\texttt{3S-like}}
\newcommand{\measplike}{\texttt{ME-ASP-like}}

\newcommand{\irace}{\textit{I/F-race}}
\newcommand{\gga}{\textit{GGA}}
\newcommand{\smac}{\textit{SMAC}}
\newcommand{\paramils}{\textit{ParamILS}}
\newcommand{\spearmint}{\textit{Spearmint}}
\newcommand{\tpe}{\textit{TPE}}


\usepackage{pifont}
\newcommand{\itarrow}{\mbox{\Pisymbol{pzd}{229}}}
\newcommand{\ithook}{\mbox{\Pisymbol{pzd}{52}}}
\newcommand{\itcross}{\mbox{\Pisymbol{pzd}{56}}}
\newcommand{\ithand}{\mbox{\raisebox{-1pt}{\Pisymbol{pzd}{43}}}}

%\DeclareMathOperator*{\argmax}{arg\,max}

\newcommand{\ie}{{\it{}i.e.\/}}
\newcommand{\eg}{{\it{}e.g.\/}}
\newcommand{\cf}{{\it{}cf.\/}}
\newcommand{\wrt}{\mbox{w.r.t.}}
\newcommand{\vs}{{\it{}vs\/}}
\newcommand{\vsp}{{\it{}vs\/}}
\newcommand{\etc}{{\copyedit{etc.}}}
\newcommand{\etal}{{\it{}et al.\/}}

\newcommand{\pscProc}{{\bf procedure}}
\newcommand{\pscBegin}{{\bf begin}}
\newcommand{\pscEnd}{{\bf end}}
\newcommand{\pscEndIf}{{\bf endif}}
\newcommand{\pscFor}{{\bf for}}
\newcommand{\pscEach}{{\bf each}}
\newcommand{\pscThen}{{\bf then}}
\newcommand{\pscElse}{{\bf else}}
\newcommand{\pscWhile}{{\bf while}}
\newcommand{\pscIf}{{\bf if}}
\newcommand{\pscRepeat}{{\bf repeat}}
\newcommand{\pscUntil}{{\bf until}}
\newcommand{\pscWithProb}{{\bf with probability}}
\newcommand{\pscOtherwise}{{\bf otherwise}}
\newcommand{\pscDo}{{\bf do}}
\newcommand{\pscTo}{{\bf to}}
\newcommand{\pscOr}{{\bf or}}
\newcommand{\pscAnd}{{\bf and}}
\newcommand{\pscNot}{{\bf not}}
\newcommand{\pscFalse}{{\bf false}}
\newcommand{\pscEachElOf}{{\bf each element of}}
\newcommand{\pscReturn}{{\bf return}}

%\newcommand{\param}[1]{{\sl{}#1}}
\newcommand{\var}[1]{{\it{}#1}}
\newcommand{\cond}[1]{{\sf{}#1}}
%\newcommand{\state}[1]{{\sf{}#1}}
%\newcommand{\func}[1]{{\sl{}#1}}
\newcommand{\set}[1]{{\Bbb #1}}
%\newcommand{\inst}[1]{{\tt{}#1}}
\newcommand{\myurl}[1]{{\small\sf #1}}

\newcommand{\Nats}{{\Bbb N}}
\newcommand{\Reals}{{\Bbb R}}
\newcommand{\extset}[2]{\{#1 \; | \; #2\}}

\newcommand{\vbar}{$\,\;|$\hspace*{-1em}\raisebox{-0.3mm}{$\,\;\;|$}}
\newcommand{\vendbar}{\raisebox{+0.4mm}{$\,\;|$}}
\newcommand{\vend}{$\,\:\lfloor$}


\newcommand{\goleft}[2][.7]{\parbox[t]{#1\linewidth}{\strut\raggedright #2\strut}}
\newcommand{\rightimage}[2][.3]{\mbox{}\hfill\raisebox{1em-\height}[0pt][0pt]{\includegraphics[width=#1\linewidth]{#2}}\vspace*{-\baselineskip}}






\newcommand{\lz}{\vspace{0.5cm}}
\newcommand{\thetab}{\bm{\weights}}
\newcommand{\zero}{\mathbf{0}}
\newcommand{\Xmat}{\mathbf{X}}
\newcommand{\Kmat}{\mathbf{K}}
\newcommand{\ydat}{\mathbf{y}}
\newcommand{\id}{\boldsymbol{I}}
\newcommand{\Amat}{\mathbf{A}}
\newcommand{\Xspace}{\mathcal{X}}                                           
\newcommand{\Yspace}{\mathcal{Y}}
\newcommand{\ls}{\ell}
\newcommand{\natnum}{\mathbb{N}}
\newcommand{\intnum}{\mathbb{Z}}

\usepackage{fontawesome}
\usepackage{dirtytalk}
\usepackage{csquotes}

%\begin{frame}[c]{}
%\centering
%\huge
%\textbf{}
%\end{frame}


%\item[\faLightbulbO]
\title[AutoML: GPs]{AutoML: Gaussian Processes} % week title
\subtitle{Mean Functions for Gaussian Processes} % video title
\author[Marius Lindauer]{\underline{Bernd Bischl} \and Frank Hutter \and Lars Kotthoff\newline \and Marius Lindauer \and Joaquin Vanschoren}
\institute{}
\date{}



% \AtBeginSection[] % Do nothing for \section*
% {
%   \begin{frame}{Outline}
%     \bigskip
%     \vfill
%     \tableofcontents[currentsection]
%   \end{frame}
% }

\begin{document}
\maketitle
	


\begin{frame}[c,allowframebreaks]{The Role of Mean Functions}

\begin{itemize}
\item It is common but by no means necessary to consider GPs with zero-mean functions: $$m(\x) \equiv 0$$

\vspace{3mm}
\item This is not necessarily a drastic limitation, since the mean of the posterior process is not confined to be zero:
\vspace{-3mm}
  $$\bm{f}_* | \Xmat_*, \Xmat, \bm{f} \sim \mathcal{N}(\Kmat_{*}^\top\Kmat^{-1}\bm{f}, \Kmat_{**} - \Kmat_*^\top \Kmat ^{-1}\Kmat_*).$$
  
\vspace{3mm}
\item There are several reasons why one might wish to explicitly model a mean function, including interpretability, convenience of expressing prior informations, etc.

%%%%%%%%%%%%%%%%%%%%%%%%%%%%%%%%%%%%%%%%%%%%%%%%%%%%%%%%%%%%%%%%%%%%%%%%%%%%%%%%%%%%

\framebreak

\item When assuming a non-zero mean GP prior $\gp (m(\x), k(\x, \x^\prime)$ with mean $m(\x)$, the predictive mean becomes:\vspace{-5mm}$$m(\Xmat_*) + \Kmat_*\Kmat_y^{-1}\left(\bm{y} - m(\Xmat)\right)$$\vspace{-7mm}but the predictive variance remains unchanged. 

\vspace{20mm}
\item Gaussian processes with non-zero mean Gaussian process priors are also called \textbf{Gaussian processes with trend}.  

\end{itemize}

\framebreak
%%%%%%%%%%%%%%%%%%%%%%%%%%%%%%%%%%%%%%%%%%%%%%%%%%%%%%%%%%%%%%%%%%%%%%%%%%%%%%%%%%%%


\begin{figure}
\includegraphics[width=0.6\textwidth]{figure_man/gp-sample-nonzero/gp-sample-nonzero-1-1.pdf}
\end{figure}

\framebreak
%%%%%%%%%%%%%%%%%%%%%%%%%%%%%%%%%%%%%%%%%%%%%%%%%%%%%%%%%%%%%%%%%%%%%%%%%%%%%%%%%%%%

\foreach \x in{1, 2, 3, 4} {
\begin{figure}
  \includegraphics[width=0.6\textwidth]{figure_man/gp-sample-nonzero/gp-sample-nonzero-2-\x.pdf}
\framebreak
\end{figure}
}

\framebreak
%%%%%%%%%%%%%%%%%%%%%%%%%%%%%%%%%%%%%%%%%%%%%%%%%%%%%%%%%%%%%%%%%%%%%%%%%%%%%%%%%%%%
\begin{itemize}
\item In practice it is often difficult to specify a fixed mean function. Instead, it may be more convenient to specify a few fixed basis functions, whose coefficients, $\bm{\beta}$, are to be inferred from the data.
\vspace{2mm}
\item Consider 
$$g(\x) = f(\x)+\bm{h}(\x)^\top\bm{\beta}, \text{   with } f(\x) \sim\gp \left(0, k(\x, \x^\prime)\right)$$ 

where $f(\x)$ is a zero mean GP, $\bm{h}(\x)$ are a set of fixed basis functions, and $\bm{\beta}$ are additional parameters. 
\vspace{2mm}
\item This formulation expresses that the data is close to a global linear model with the residuals being modelled by a GP.
\vspace{2mm}
\item By taking the prior on $\bm{\beta}$ to be Gaussian, $\bm{\beta} \sim \mathcal{N}(\mathbf{b}, B)$, another GP with an added contribution in the covariance function can be obtained: 
\vspace{-2mm}
$$g(\x) \sim \gp\left(\mathbf{h}(\x^\top \mathbf{b}, k\left(\x, \x^\prime\right)+\mathbf{h}(\x)^\top B\,\, \mathbf{h}\left(\x^\prime\right)\right)$$

\end{itemize}


\framebreak
%%%%%%%%%%%%%%%%%%%%%%%%%%%%%%%%%%%%%%%%%%%%%%%%%%%%%%%%%%%%%%%%%%%%%%%%%%%%%%%%%%%%
\begin{itemize}

\item The mean and covariance functions of $g(\x)$ will be then:
\vspace{-2mm}

$$\begin{aligned}
\overline{\mathbf{g}}\left(X_{*}\right) &=\overline{\mathbf{f}}\left(X_{*}\right)+R^{\top} \overline{\boldsymbol{\beta}}, \\
\operatorname{cov}\left(\mathbf{g}_{*}\right) &=\operatorname{cov}\left(\mathbf{f}_{*}\right)+R^{\top}(B^{-1}+H K_{y}^{-1} H^{\top})^{-1} R,
\end{aligned}\label{41}$$ where $\overline{\boldsymbol{\beta}}=(B^{-1}+H K_{y}^{-1} H^{\top})^{-1}(H K_{y}^{-1} \mathbf{y}+B^{-1} \mathbf{b})$ and  $R=H_{*}-H K_{y}^{-1} K_{*}.$

\vspace{4mm}
\item Notice that $\overline{\boldsymbol{\beta}}$ is the mean of the global linear model parameters and the $H$ matrix collects the $\mathbf{h}(\x)$ vectors for all training and test cases.

\vspace{4mm}
\item Hence, the predictive mean is the sum of the mean linear output and what the GP model predicts from the residuals. The covariance is the usual covariance term plus a non negative contribution.
\end{itemize}

\framebreak
%%%%%%%%%%%%%%%%%%%%%%%%%%%%%%%%%%%%%%%%%%%%%%%%%%%%%%%%%%%%%%%%%%%%%%%%%%%%%%%%%%%%

\begin{itemize}

\item In the limiting case where the prior on the $\boldsymbol{\beta}$ becomes vague ($B^{-1} \rightarrow O$), the predictive distribution becomes independent of $\boldsymbol{b}$:
\vspace{-2mm}

$$\begin{aligned} 
\overline{\mathbf{g}}\left(X_{*}\right) &=\overline{\mathbf{f}}\left(X_{*}\right)+R^{\top} \overline{\boldsymbol{\beta}}, \\
\operatorname{cov}\left(\mathbf{g}_{*}\right) &=\operatorname{cov}(\mathbf{f}_{*})+R^{\top}(H K_{y}^{-1} H^{\top})^{-1} R, 
\end{aligned}
$$ where the limiting $\overline{\boldsymbol{\beta}}=(H K_{y}^{-1} H^{\top})^{-1} H K_{y}^{-1} \mathbf{y}$.

\vspace{1cm}

\item To make predictions under the limit $B^{-1} \rightarrow O$, it is important to use the above equation. Otherwise, by na{\"\i}vely plugging the modified covariance function into the standard prediction equations, the entries of the covariance function tend to infinity making it unsuitable for numerical implementation.

\end{itemize}

\framebreak
%%%%%%%%%%%%%%%%%%%%%%%%%%%%%%%%%%%%%%%%%%%%%%%%%%%%%%%%%%%%%%%%%%%%%%%%%%%%%%%%%%%%
\begin{itemize}

\item The marginal likelihood for the model with a Gaussian prior $\bm{\beta} \sim \mathcal{N}(\mathbf{b}, B)$ can be expressed as what follows (the explicit mean has been included).

%\begin{footnotesize}
$$\begin{aligned} 
\log p(\mathbf{y} | X, \mathbf{b}, B)=&-\frac{1}{2}\left(H^{\top} \mathbf{b}-\mathbf{y}\right)^{\top}\left(K_{y}+H^{\top} B H\right)^{-1}\left(H^{\top} \mathbf{b}-\mathbf{y}\right) \\ &-\frac{1}{2} \log \left|K_{y}+H^{\top} B H\right|-\frac{n}{2} \log 2 \pi. 
\end{aligned}$$
%\end{footnotesize}

\lz
\lz

\item We are interested in exploring the limit where $B^{-1} \rightarrow O$, i.e. when the prior is vague. In this limit the mean of the prior is irrelevant, so without loss of generality we assume for now that the mean is zero, $\mathbf{b} = \mathbf{0}$.

\end{itemize}

\framebreak
%%%%%%%%%%%%%%%%%%%%%%%%%%%%%%%%%%%%%%%%%%%%%%%%%%%%%%%%%%%%%%%%%%%%%%%%%%%%%%%%%%%%
\begin{itemize}

\item This assumption gives: 
\vspace{-5mm}
%\begin{footnotesize} 
$$\begin{aligned} 
\log p(\mathbf{y} | X, \mathbf{b}=\mathbf{0}, B)=&-\frac{1}{2} \mathbf{y}^{\top} K_{y}^{-1} \mathbf{y}+\frac{1}{2} \mathbf{y}^{\top} C \mathbf{y} \\ &-\frac{1}{2} \log \left|K_{y}\right|-\frac{1}{2} \log |B|-\frac{1}{2} \log |A|-\frac{n}{2} \log 2 \pi, 
\end{aligned}$$
where $A=B^{-1}+H K_{y}^{-1} H^{\top}$ and $C=K_{y}^{-1} H^{\top} A^{-1} H K_{y}^{-1}$.
%\end{footnotesize}

\vspace{4mm}
\item The behaviour of the log marginal likelihood in the limiting case can be explored now. The variances of the Gaussian in the directions spanned by columns of $H^\top$ will become infinite, and it will require special treatment. 

\vspace{4mm}
\item[\faLightbulbO] The log marginal likelihood consists of three terms: a quadratic form in $\mathbf{y}$, a log determinant term, and a term involving $\log 2 \pi$.

\end{itemize}

\framebreak
%%%%%%%%%%%%%%%%%%%%%%%%%%%%%%%%%%%%%%%%%%%%%%%%%%%%%%%%%%%%%%%%%%%%%%%%%%%%%%%%%%%%
\begin{itemize}

\item Performing an eigendecomposition of the covariance matrix, the contributions to quadratic form term from the infinite-variance directions will be zero. However, the log determinant term will tend to minus infinity. 

\vspace{3mm}
\item The standard solution in this case is to project $\mathbf{y}$ onto the directions orthogonal to the span of $H^\top$ and compute the marginal likelihood in this subspace. 

\vspace{3mm}
\item By taking $m$ to denote the rank of $H^\top$, we can discard the terms $-\frac{1}{2} \log |B|-\frac{m}{2} \log 2 \pi$ from the previous equation to give:
\vspace{-3mm}
%\begin{footnotesize}
$$\log p(\mathbf{y} | X)=-\frac{1}{2} \mathbf{y}^{\top} K_{y}^{-1} \mathbf{y}+\frac{1}{2} \mathbf{y}^{\top} C \mathbf{y}-\frac{1}{2} \log \left|K_{y}\right|-\frac{1}{2} \log |A|-\frac{n-m}{2} \log 2 \pi,$$
where $A=H K_{y}^{-1} H^{\top}$ and $C=K_{y}^{-1} H^{\top} A^{-1} H K_{y}^{-1}$.
%\end{footnotesize}

\end{itemize}

\framebreak
%%%%%%%%%%%%%%%%%%%%%%%%%%%%%%%%%%%%%%%%%%%%%%%%%%%%%%%%%%%%%%%%%%%%%%%%%%%%%%%%%%%%



\framebreak
%%%%%%%%%%%%%%%%%%%%%%%%%%%%%%%%%%%%%%%%%%%%%%%%%%%%%%%%%%%%%%%%%%%%%%%%%%%%%%%%%%%%



\framebreak
%%%%%%%%%%%%%%%%%%%%%%%%%%%%%%%%%%%%%%%%%%%%%%%%%%%%%%%%%%%%%%%%%%%%%%%%%%%%%%%%%%%%



\framebreak
%%%%%%%%%%%%%%%%%%%%%%%%%%%%%%%%%%%%%%%%%%%%%%%%%%%%%%%%%%%%%%%%%%%%%%%%%%%%%%%%%%%%



\framebreak
%%%%%%%%%%%%%%%%%%%%%%%%%%%%%%%%%%%%%%%%%%%%%%%%%%%%%%%%%%%%%%%%%%%%%%%%%%%%%%%%%%%%



\framebreak
%%%%%%%%%%%%%%%%%%%%%%%%%%%%%%%%%%%%%%%%%%%%%%%%%%%%%%%%%%%%%%%%%%%%%%%%%%%%%%%%%%%%



\framebreak
%%%%%%%%%%%%%%%%%%%%%%%%%%%%%%%%%%%%%%%%%%%%%%%%%%%%%%%%%%%%%%%%%%%%%%%%%%%%%%%%%%%%



\framebreak
%%%%%%%%%%%%%%%%%%%%%%%%%%%%%%%%%%%%%%%%%%%%%%%%%%%%%%%%%%%%%%%%%%%%%%%%%%%%%%%%%%%%



\framebreak
%%%%%%%%%%%%%%%%%%%%%%%%%%%%%%%%%%%%%%%%%%%%%%%%%%%%%%%%%%%%%%%%%%%%%%%%%%%%%%%%%%%%



\framebreak
%%%%%%%%%%%%%%%%%%%%%%%%%%%%%%%%%%%%%%%%%%%%%%%%%%%%%%%%%%%%%%%%%%%%%%%%%%%%%%%%%%%%




\end{frame}
%%%%%%%%%%%%%%%%%%%%%%%%%%%%%%%%%%%%%%%%%%%%%%%%%%%%%%%%%%%%%%%%%%%%%%%%%%%%%%%%%%%%


\end{document}
